\input{../tex-inputs/latex-leseansicht-vorspann}

\section[Arthur Schnitzler an Gustav Schwarzkopf, 18. 8. 1903]{L04440 Arthur Schnitzler an Gustav Schwarzkopf, 18. 8. 1903}
\nopagebreak\mylabel{L04440v}
\rehead{ }\normalsize\beginnumbering\briefempfaengerindex{Schwarzkopf, Gustav@\textsc{Schwarzkopf, Gustav}!zzzSchnitzler, Arthur@\emph{von Arthur Schnitzler}!1903-08-181@{18. 8. 1903}|(be}
\toendnotes[C]{\smallbreak\pagebreak[2]}
\correspDesc{Versand  durch Arthur Schnitzler am 18. 8. 1903 in Madonna di Campiglio
\newline{}Übermittlung  in Madonna di Campiglio
\newline{}Zustellung  am 20. 8. 1903 in Wien
\newline{}Erhalt  durch Gustav Schwarzkopf am 20. 8. 1903 in Wien}\toendnotes[C]{\smallbreak}
\Standort{DLA, A:Schnitzler, HS.1985.1.1897.}
\physDesc{Bildpostkarte, 76 Zeichen
\newline{}Handschrift: Bleistift, deutsche Kurrent
\newline{}Versand: 1) Stempel: »\nobreak{}\oindex{Madonna di Campiglio@\textbf{Madonna di Campiglio}|pwk}\textcolor{gray}{M}adonna\nobreak{}«.   2) Stempel: »\nobreak{}\oindex{I., Innere Stadt@\textbf{I., Innere Stadt}|pwk}Wien 1, 20. \textcolor{gray}{8}. 0\textcolor{gray}{3}, 8–9½, \textcolor{gray}{Bestellt}\nobreak{}«. }\pstart{}{\pb}Herrn Guſtav Schwarzkopf\pend{}\pstart{}Wien I\oindex{I., Innere Stadt@\textbf{I., Innere Stadt}|pw}.\pend{}\pstart{}Tiefer Graben 23\oindex{Wien@\textbf{Wien}!I., Innere Stadt@\textbf{I., Innere Stadt}!Tiefer Graben 23@\textbf{Tiefer Graben 23}|pw}.\pend{}{\bigskip}
\pstart
           {\pb}\textcolor{gray}{\textbf{Madonna di Campiglio}}\oindex{Madonna di Campiglio@\textbf{Madonna di Campiglio}|pw}\pend
           \vspace{1em}
\pstart
           {\pb}18. 8. 903\pend
           \vspace{0.5em}
\pstart
           Herzlichen Gruſs\pend
           \pstart \spacefill\mbox{A.}\pend{}\selectlanguage{ngerman}\endnumbering\briefempfaengerindex{Schwarzkopf, Gustav@\textsc{Schwarzkopf, Gustav}!zzzSchnitzler, Arthur@\emph{von Arthur Schnitzler}!1903-08-181@{18. 8. 1903}|)be}\mylabel{L04440h}
\begin{anhang}
\end{anhang}\newcommand{\dateiname}{L04440}\newcommand{\titel}{Arthur Schnitzler an Gustav Schwarzkopf, 18. 8. 1903}\newcommand{\editorInnen}{Herausgegeben von Jahnke, SelmaMüller, Martin Anton}\input{../tex-inputs/latex-leseansicht-abspann}
